\section{Infirmity, Death, and the Making of a Doctor (604 and After)}

\subsection{The final years and death (12 March 604)}

Gregory governed the Church from a body that was steadily failing. Already in
598 he wrote to Eulogius of Alexandria that the messenger ``found me sick, and
has left me sick'' (\work{Ep.}\ VIII.30). By the last winter of his life,
severe gout confined him to his bed for long stretches and sometimes left him
scarcely able to speak. Yet the seventeen surviving letters of Book XIV
(Indiction VII, September 603 to March 604, summarized in Barmby's
Prolegomena) show no slackening of his vigilance over churches, monasteries,
and peace in Italy.

His final surviving letter to the Lombard Queen Theodelinda (\work{Ep.}\
XIV.12, reported in Barmby, Prolegomena) forms a fitting testament. Writing
``between life and death'' while still speaking of the remote hope of recovery,
he congratulated the queen on the Catholic baptism of her infant son Adaloald
(heir to the Arian King Agilulf), sent a cross containing a relic of the True
Cross for the boy and jeweled rings for his sister, urged her to thank Agilulf
for the renewed truce of September 603, and promised to answer her remaining
objections concerning the Fifth Council if God granted him strength. He never
had the chance.

Pope Gregory died at Rome on 12 March 604, after thirteen years and six months
in the chair of Saint Peter. He was buried in the portico of Old Saint Peter's
Basilica. His epitaph---preserved in Bede (\work{Historia Ecclesiastica} II.1,
reported in Barmby's Prolegomena)---saluted him with a title that fused his two
worlds into two Latin words: \latin{Consul Dei}, ``God's Consul.'' He had
turned the old Roman magistracy inside out, feeding the hungry, ransoming the
captive, and defending the city not with legions, but with charity and the
Gospel.

\subsection{Burial, earliest cult, and veneration}

Remarkably, Gregory's formal literary cult bloomed first in the northern island
he had never seen. While seventh-century Rome faced famine, political tension,
and resentment from some clerical factions who had chafed under a monastic
pope, Anglo-Saxon England revered him immediately as its spiritual father. At
the Northumbrian abbey of Whitby between roughly 704 and 714, an anonymous
monk composed the earliest extant \work{Vita Gregorii}. Twenty years later,
Bede dedicated the opening chapter of Book II of his \work{Historia
Ecclesiastica} (c.\ 731) to Gregory, calling him ``our apostle'' because
through his zeal the English had passed from idols to Christ.

In Rome itself, devotion took firm liturgical root during the seventh and
eighth centuries. By the time Pope Hadrian I sent the Roman Sacramentary to
Charlemagne (784--791), Gregory's feast on 12 March (\latin{Natale Sancti
Gregorii Papae}) already stood in the papal calendar with its own proper
prayers---asking that ``we who are weighed down by the burden of our sins may
be lifted up by his prayers before thee'' (Wilson 1915, p.\ 25). In the ninth
century, Pope Gregory IV translated his remains from the portico into the
interior of Saint Peter's Basilica (his relics now rest beneath the Altar of
Saint Gregory in the Clementine Chapel of New Saint Peter's). The 1902
\work{Martyrologium Romanum} records two distinct Roman feasts: on 12 March,
his \latin{dies natalis}, honoring him as ``Pope and eminent Doctor of the
Church, who for his illustrious deeds and for converting the English to the
faith of Christ was called the Great and named Apostle of the English''
(\latin{Magnus est dictus, et Anglorum Apostolus appellatus}); and on 3
September, the commemoration of his episcopal ordination, ``who, compelled to
undergo that burden, shone forth from a loftier throne with brighter rays of
sanctity throughout the world'' (\work{Martyrologium Romanum}, 1902, pp.\ 36,
132). Following the 1969 revision of the General Roman Calendar, the universal
Roman memorial moved from 12 March (which always falls in Lent) to 3 September,
the anniversary of his consecration.

\subsection{Doctor of the Church and Western inheritance}

Throughout the Middle Ages, Gregory stood beside Ambrose, Augustine, and
Jerome as one of the four pillars of Latin Christian teaching. In 1298, Pope
Boniface VIII formally confirmed that universal custom by ordering the feasts
of the four great Western Doctors kept with double rank throughout the Church.
In the Byzantine East as well, Gregory is venerated as Saint Gregory
the Dialogist (\latin{Dialogos}), whose name became attached in Eastern
liturgical usage to the Lenten Liturgy of the Presanctified Gifts.

His theological influence ran less through speculative originality than
through synthesis, pastoral wisdom, and moral clarity. Medieval monastic
libraries treated the \work{Moralia in Iob} as an inexhaustible quarry for
contemplative life and the taxonomy of the seven capital vices; bishops studied
the \work{Regula Pastoralis} as their mirror of office; scholastic masters,
including Saint Thomas Aquinas throughout the \work{Summa Theologiae}, cited
Gregory constantly on angels, contemplation, grace, penance, and the active
life. Modern historians have stripped away the legendary varnish only to
confirm the medieval instinct: standing directly on the fault line where the
Roman Empire dissolved into early medieval Europe, Gregory carried the Gospel,
Latin patristic learning, Benedictine monasticism, and Roman pastoral order
safely across the divide.
